\documentclass[12pt,letterpaper]{article}
\usepackage[utf8]{inputenc}
\usepackage[spanish,es-noshorthands]{babel}
\usepackage{amsmath}
\usepackage{amsfonts}
\usepackage{amssymb}
\usepackage{graphicx}
\usepackage{hyperref}
\usepackage{geometry}
\usepackage{xcolor}
\usepackage{tcolorbox}
\usepackage{booktabs}

\geometry{left=3cm,right=3cm,top=2.5cm,bottom=2.5cm}

\title{
    \vspace{-2cm}
    \begin{center}
        \Large \textbf{Universidad de Oriente}\\
        \large \textbf{Integración de Inteligencia Artificial en el Flujo de Trabajo Académico y de Diseño Electrónico}
    \end{center}
    \vspace{1cm}
}
\author{Informe}
\date{\today}

\begin{document}

\maketitle

\vspace{-1cm}
\begin{flushright}
\textbf{Para:} Prof. Manuel Centeno\\
\textit{Director del Equipo de Inteligencia Artificial}\\
\textit{Departamento de Matemáticas, Universidad de Oriente}
\end{flushright}
\vspace{0.5cm}

\section*{Estimado Profesor Centeno,}

Es un placer saludarle. Conociendo el excelente trabajo que lidera en el área de Inteligencia Artificial dentro del Departamento de Matemáticas de nuestra alma máter, me permito compartir con usted los avances y la arquitectura del sistema que hemos desarrollado e implementado recientemente en el espacio de trabajo multidisciplinario de \textbf{Electrónica}. 

El objetivo de este proyecto ha sido potenciar la eficiencia en la docencia, la redacción académica y el Diseño de Circuitos Integrados (EDA), apalancando Modelos de Lenguaje Grande (LLMs) de código abierto y APIs multimodales bajo una arquitectura rigurosamente estructurada.

\section{Arquitectura del Sistema: Diseño de 3 Capas}

Para evitar la fragilidad común en los agentes IA que mezclan lógica probabilística con ejecución determinista, hemos implementado un marco estricto de \textbf{3 capas} inspirado en las mejores prácticas de ingeniería de software para IA:

\begin{enumerate}
    \item \textbf{Capa 1 (Directives):} Archivos YAML que contienen los Procedimientos Operativos Estándar (SOPs). Definen las reglas del negocio, rúbricas de evaluación y configuraciones estáticas sin incluir lógica de programación.
    \item \textbf{Capa 2 (Orchestration):} Scripts en Python (Orquestadores) y el Agente IA en sí. Se encargan del enrutamiento, la toma de decisiones basada en el contexto y la interacción con los LLMs.
    \item \textbf{Capa 3 (Execution):} Scripts deterministas, con una única responsabilidad, que ejecutan acciones como compilar código, consumir APIs, gestionar archivos, etc.
\end{enumerate}

\section{Componentes Tecnológicos Destacados}

El ecosistema está construido principalmente en \textbf{Python 3} y \textbf{\LaTeX}, utilizando el framework \textbf{LangChain} para la orquestación semántica. Entre los módulos más relevantes destacan:

\subsection{Sistema RAG Académico (Retrieval-Augmented Generation)}
Hemos construido una base de datos vectorial local utilizando \textbf{ChromaDB} y modelos de embeddings multilingües (\texttt{paraphrase-multilingual-MiniLM}). Esto permite a nuestro Agente (alimentado por Llama 3 vía la API de Groq) consultar instantáneamente miles de páginas de documentos \texttt{.tex}, \texttt{.md} y \texttt{.pdf} de nuestros cursos de Electrónica, Física y planes de estudio.

\subsection{Agente EDA (Diseño Electrónico Automatizado)}
Desarrollamos un agente capaz de analizar diagramas de circuitos escritos nativamente en \LaTeX\ (usando el paquete \texttt{circuitikz}), extraer las listas de conexiones (netlists) y la lista de materiales (BOM), para luego generar automáticamente los archivos JSON compatibles con herramientas estándar de la industria como \textbf{EasyEDA}.

\subsection{Evaluación Multimodal de Exámenes y Prácticas}
Mediante el uso de modelos multimodales (Gemini / OpenRouter), el sistema puede leer escaneos en PDF de los exámenes de los estudiantes o de sus prácticas de laboratorio, evaluar su desempeño comparándolo contra rúbricas estrictas definidas en YAML (Capa 1), y generar de forma autónoma un informe de retroalimentación detallado, compilado directamente en un archivo PDF mediante \LaTeX.

\subsection{Protocolo de Contexto de Modelos (MCP) y Telegram Gateway}
Una de las implementaciones más recientes y robustas ha sido la estandarización de nuestras herramientas a través de servidores \textbf{MCP (Model Context Protocol)}. Para interactuar con estos servidores de forma remota (desde un teléfono móvil, por ejemplo), construimos un \textbf{Gateway mediante un Bot de Telegram}.

Lo destacable de este Gateway es su enfoque de seguridad:
\begin{itemize}
    \item Utiliza el mecanismo de \textbf{Long Polling} (conexiones salientes), lo que significa que el sistema opera de forma segura detrás del firewall de la red local sin requerir redirección de puertos ni túneles expuestos a internet (como ngrok).
    \item Mantiene una lista blanca estricta de usuarios autorizados.
    \item El bot actúa como un cliente MCP local, tomando los comandos del usuario y enrutándolos hacia servidores especializados aislados (compilador LaTeX, diagnóstico de hardware, evaluadores).
\end{itemize}

\section{Perspectivas a Futuro}

El cruce entre la Inteligencia Artificial y nuestras disciplinas técnicas abre un abanico de posibilidades inmenso. El uso de \LaTeX\ como lenguaje puente entre la matemática/electrónica y el LLM ha demostrado ser sumamente eficaz, reduciendo las alucinaciones del modelo en formulaciones complejas.

Como siguiente fase evolutiva de este ecosistema, y sujeto a la disponibilidad de recursos, hemos estructurado el diseño formal de un \textbf{Clúster Beowulf para Cómputo en Paralelo}. Esta infraestructura aprovechará hardware preexistente (como un servidor Dell PowerEdge R610 y estaciones de trabajo) interconectado mediante una topología en estrella Gigabit. Su propósito será acelerar las cargas de trabajo más demandantes de nuestro orquestador, como el entrenamiento local e inferencia de Modelos de Lenguaje (LLMs) y la simulación física de alta fidelidad, democratizando así el acceso a la Computación de Alto Rendimiento (HPC) dentro de nuestro departamento y \textbf{eliminando por completo la necesidad de adquirir costosas tarjetas gráficas (GPUs)} gracias a la distribución del cómputo directamente sobre los procesadores (CPUs) del clúster.

Me encantaría tener la oportunidad de mostrarle este sistema en funcionamiento, intercambiar ideas y, si lo considera pertinente, explorar posibles sinergias o líneas de investigación conjuntas con el equipo de Inteligencia Artificial que usted dirige.

Quedo a su entera disposición.

\vspace{1cm}
Atentamente,

\vspace{1.5cm}
\noindent\rule{6cm}{0.4pt}\\
\textbf{Prof. César Rodríguez}\\
Dpto. Física, Universidad de Oriente-NS\\
Espacio de Trabajo: ELECTRONICA

\end{document}
